\documentclass[11pt]{article}
\usepackage[spanish,es-noshorthands,es-tabla]{babel}
\newcommand{\imgdir}{images/es}
% Shared layout for the Illuminati manual. Each manual-<lang>.tex loads this
% after \documentclass and sets \manuallang (babel name) and \imgdir.
% Build with XeLaTeX (see build.ps1); the fonts are the app's own.

\usepackage{fontspec}
\setmainfont{Literata}[
  Path=../../assets/fonts/literata/, Extension=.ttf,
  UprightFont=*-Regular, BoldFont=*-Bold, ItalicFont=*-Italic, BoldItalicFont=*-BoldItalic]
\setsansfont{AtkinsonHyperlegible}[
  Path=../../assets/fonts/atkinson/, Extension=.ttf,
  UprightFont=*-Regular, BoldFont=*-Bold, ItalicFont=*-Italic, BoldItalicFont=*-BoldItalic]
\newfontfamily\dyslexicfont{OpenDyslexic}[
  Path=../../assets/fonts/opendyslexic/, Extension=.otf,
  UprightFont=*-Regular, BoldFont=*-Bold, ItalicFont=*-Italic, BoldItalicFont=*-Bold-Italic]
\newfontfamily\garamondfont{EBGaramond}[
  Path=../../assets/fonts/ebgaramond/, Extension=.ttf,
  UprightFont=*-Regular, BoldFont=*-Bold, ItalicFont=*-Italic, BoldItalicFont=*-BoldItalic]

\usepackage[a4paper, margin=2.4cm, headheight=14pt]{geometry}
\usepackage{graphicx}
\usepackage[table]{xcolor}
\usepackage{booktabs, tabularx, array}
\usepackage{enumitem}
\usepackage{fancyhdr}
\usepackage{titlesec}
\usepackage[most]{tcolorbox}
\usepackage[hidelinks]{hyperref}

% Colours from the app's logo.
\definecolor{violet}{HTML}{34206A}
\definecolor{violetlight}{HTML}{EDE8F7}
\definecolor{gold}{HTML}{C9A227}
\definecolor{goldlight}{HTML}{FBF4DF}

\setlength{\parskip}{0.55em}
\setlength{\parindent}{0pt}
\setlist{itemsep=0.25em, topsep=0.3em}
\linespread{1.08}

\titleformat{\section}{\sffamily\Large\bfseries\color{violet}}{\thesection}{0.8em}{}
\titleformat{\subsection}{\sffamily\large\bfseries\color{violet}}{\thesubsection}{0.7em}{}
\titlespacing*{\section}{0pt}{1.6em}{0.6em}
\titlespacing*{\subsection}{0pt}{1.1em}{0.35em}
% Don't leave a heading alone at the bottom of a page.
\usepackage{needspace, etoolbox}
\pretocmd{\section}{\needspace{8\baselineskip}}{}{}
\pretocmd{\subsection}{\needspace{5\baselineskip}}{}{}

\pagestyle{fancy}
\fancyhf{}
\fancyhead[L]{\sffamily\small\color{violet} Illuminati}
\fancyhead[R]{\sffamily\small\color{gray}\nouppercase{\rightmark}}
\fancyfoot[C]{\sffamily\small\thepage}
\renewcommand{\headrulewidth}{0.3pt}
% \rightmark is the first section mark on the page.
\renewcommand{\sectionmark}[1]{\markboth{#1}{#1}}
\renewcommand{\subsectionmark}[1]{}

% Keep catalogue names whole in every language.
\babelhyphenation{Wiki-source Guten-berg}

% A key or key combination, e.g. \kbd{Ctrl+F}.
\newcommand{\kbd}[1]{\tikz[baseline=(k.base)]\node[draw=gray!60, fill=gray!8,
  rounded corners=2pt, inner sep=2pt, font=\sffamily\small](k){#1};}
\usepackage{tikz}

% Arrow keys (the sans font has no arrow characters).
\newcommand{\kright}{\kbd{\ensuremath{\rightarrow}}}
\newcommand{\kleft}{\kbd{\ensuremath{\leftarrow}}}
\newcommand{\kup}{\kbd{\ensuremath{\uparrow}}}
\newcommand{\kdown}{\kbd{\ensuremath{\downarrow}}}

% A button, menu item or label as it appears in the app.
\newcommand{\ui}[1]{\textsf{\textbf{#1}}}

% Screenshot with caption; skipped if the image file isn't there.
\newcommand{\shot}[2]{%
  \IfFileExists{\imgdir/#1.jpg}{%
    \begin{figure}[htbp]\centering
      \fcolorbox{gray!40}{white}{\includegraphics[width=0.96\textwidth]{\imgdir/#1.jpg}}
      \caption{#2}
    \end{figure}}{}}

\newtcolorbox{tip}[1]{enhanced, colback=goldlight, colframe=gold,
  boxrule=0.6pt, arc=3pt, left=8pt, right=8pt, top=5pt, bottom=5pt,
  fonttitle=\sffamily\bfseries, coltitle=black, colbacktitle=gold!35,
  title={#1}}
\newtcolorbox{note}[1]{enhanced, colback=violetlight, colframe=violet,
  boxrule=0.6pt, arc=3pt, left=8pt, right=8pt, top=5pt, bottom=5pt,
  fonttitle=\sffamily\bfseries, coltitle=white, colbacktitle=violet,
  title={#1}}

% Two-column tables that fill the text width.
\newenvironment{keytable}[2]{%
  \renewcommand{\arraystretch}{1.3}%
  \rowcolors{2}{white}{violetlight}%
  \tabularx{\textwidth}{@{}>{\raggedright\arraybackslash}p{0.33\textwidth}X@{}}
  \toprule \textsf{\textbf{#1}} & \textsf{\textbf{#2}} \\ \midrule}{%
  \bottomrule\endtabularx}

% Title page: \maketitlepage{subtitle}{version line}{blurb}
\newcommand{\maketitlepage}[3]{%
  \begin{titlepage}
    \centering
    \vspace*{3cm}
    \includegraphics[width=4.2cm]{../../assets/icon/app_icon.png}\par\vspace{1.2cm}
    {\sffamily\fontsize{40}{46}\selectfont\bfseries\color{violet} Illuminati\par}
    \vspace{0.5cm}
    {\sffamily\LARGE #1\par}
    \vspace{1.5cm}
    {\large\itshape #3\par}
    \vfill
    {\sffamily\color{gray} #2\par}
  \end{titlepage}}

\hypersetup{pdftitle={Illuminati — Manual de usuario}, pdfauthor={Illuminati}, pdflang={es}}

\begin{document}

\maketitlepage{Manual de usuario}{Versión 1.0 · Windows, macOS, Linux y Android}%
  {Busca, descarga y lee libros que nunca tuvieron derechos de autor\\o cuyos derechos ya expiraron.}

{\setcounter{tocdepth}{1}\tableofcontents}
\newpage

% ---------------------------------------------------------------------------
\section{Bienvenida}

Illuminati es una aplicación de lectura para Windows, macOS, Linux y Android. Funciona como las
aplicaciones de libros electrónicos que quizá ya conozcas, con una diferencia: todos los libros
que ofrece son gratuitos y legales de descargar, porque sus derechos de autor expiraron o nunca
los tuvieron.

\begin{itemize}
  \item Más de 75.000 libros de \textbf{Project Gutenberg}, además de \textbf{Standard Ebooks},
        \textbf{Wikisource} e \textbf{Internet Archive}.
  \item Un lector cómodo con subrayados, notas, marcadores, búsqueda, diccionario y traducción.
  \item Tu propia biblioteca, con estantes y listas de lectura.
  \item Ajustes de pantalla y accesibilidad, incluida una tipografía pensada para la dislexia.
  \item Sin cuenta, sin suscripción y sin publicidad. Todo se guarda en tu dispositivo.
\end{itemize}

La interfaz está disponible en español, inglés, francés, alemán, italiano y portugués.

% ---------------------------------------------------------------------------
\section{Primeros pasos}

\subsection{Instalación}
\begin{itemize}
  \item \textbf{Windows:} copia la carpeta \texttt{Illuminati} completa a tu computadora y ejecuta
        \texttt{illuminati.exe}. Deja los demás archivos junto a él; el programa los necesita.
  \item \textbf{Android:} instala el archivo \texttt{.apk}. Puede que tengas que permitir la
        instalación de aplicaciones de este origen en los ajustes del teléfono.
  \item \textbf{macOS y Linux:} instala la versión preparada para tu sistema.
\end{itemize}

\subsection{La pantalla de inicio}
En una computadora, la aplicación se abre maximizada. La pantalla de inicio tiene dos partes:

\begin{enumerate}
  \item \ui{Tu biblioteca} (arriba): los libros que descargaste, empezando por el último que
        leíste, con una barra de progreso. El botón \ui{Biblioteca y estantes} abre toda tu
        biblioteca.
  \item \ui{Descubre libros gratis} (abajo): un cuadro de búsqueda, el idioma del catálogo, las
        categorías, una fila para cada catálogo adicional y, al final, el catálogo de Project
        Gutenberg.
\end{enumerate}

\shot{home}{La pantalla de inicio.}

Los dos botones de arriba a la derecha abren \ui{Pantalla y accesibilidad}
(sección~\ref{sec:display}) y \ui{Acerca de}, que indica de dónde vienen los libros y las
licencias de las tipografías.

% ---------------------------------------------------------------------------
\section{Buscar libros}

\subsection{Búsqueda}
Escribe en \ui{Buscar por título o autor}. La búsqueda:
\begin{itemize}
  \item encuentra cualquier parte de un título o de un autor: \emph{miserab} encuentra
        \emph{Los miserables};
  \item no distingue acentos ni mayúsculas: \emph{misérable} y \emph{MISERABLE} dan los mismos
        resultados;
  \item busca en todos los catálogos a la vez; cada uno muestra sus resultados en su propia fila.
\end{itemize}
Borra el cuadro con el botón \ui{×} para volver a los libros populares.

\subsection{Categorías e idioma}
Las etiquetas debajo del cuadro de búsqueda (\ui{Ficción}, \ui{Misterio}, \ui{Ciencia ficción},
\ui{Poesía}…) limitan todos los catálogos a un tema. El menú de idioma de la derecha muestra solo
los libros de un idioma, o \ui{Todos los idiomas}. Los catálogos que no tienen nada para el idioma
o la categoría elegidos se ocultan.

\subsection{Los catálogos}
\begin{keytable}{Catálogo}{Qué ofrece}
  Project Gutenberg & La mayor biblioteca gratuita, con más de 75.000 libros ordenados por
                      popularidad. Muchos tienen una edición ilustrada. \\
  Standard Ebooks   & Clásicos revisados con esmero y con una tipografía cuidada, en inglés. \\
  Wikisource        & Textos transcritos por voluntarios, en muchos idiomas. El libro electrónico
                      se prepara al descargarlo; las obras muy largas (por ejemplo, novelas en
                      varios tomos) pueden no estar disponibles. \\
  Internet Archive  & Escaneos de ediciones impresas antiguas, publicadas hace 95 años o más. El
                      texto se obtuvo automáticamente de los escaneos y puede contener errores. \\
\end{keytable}

\subsection{La página de un libro}
Toca una portada para abrir la página del libro. Muestra el resumen, los temas, el idioma, el
catálogo y una nota sobre los derechos de autor. Desde ahí puedes:
\begin{itemize}
  \item \ui{Descargar gratis}: el libro se guarda en tu dispositivo, listo para leer sin conexión.
        Cada portada del catálogo también tiene un pequeño botón de descarga.
  \item \ui{Leer ahora} o \ui{Seguir leyendo}, una vez descargado.
  \item \ui{Quiero leerlo}, \ui{Terminado} y \ui{Añadir a un estante}, incluso antes de
        descargarlo (sección~\ref{sec:library}).
  \item \ui{Obtener la edición ilustrada}, para libros de Gutenberg que descargaste sin imágenes.
\end{itemize}

\shot{details}{La página de un libro.}

% ---------------------------------------------------------------------------
\section{Tu biblioteca}\label{sec:library}

\ui{Biblioteca y estantes}, en la pantalla de inicio, abre todos los libros que guardaste: los
descargados y los que guardaste para más tarde sin descargarlos (con la portada atenuada y una
nube).

\shot{library}{La biblioteca, con sus vistas y estantes.}

\subsection{Vistas, filtro y orden}
Las etiquetas de arriba muestran una vista, cada una con su número de libros: \ui{Todos},
\ui{Descargados}, \ui{Quiero leerlo}, \ui{Leyendo}, \ui{Terminado} y, después, cada uno de tus
estantes. El cuadro de encima filtra por título o autor. El botón de ordenar (arriba a la derecha)
ordena los libros por \ui{Leídos recientemente}, \ui{Título}, \ui{Autor} (por apellido),
\ui{Progreso} o \ui{Añadidos recientemente}. La aplicación recuerda tu elección.

\subsection{Listas de lectura}
\begin{itemize}
  \item \ui{Quiero leerlo} guarda un libro para más tarde, sin descargarlo. Al abrir el libro, sale
        de esta lista.
  \item \ui{Terminado} se marca solo cuando llegas al final de un libro (puedes deshacerlo), o a
        mano.
\end{itemize}

\subsection{Estantes}
Los estantes son tus propias colecciones, como \emph{Favoritos} o \emph{Filosofía}. Un libro puede
estar en varios estantes.
\begin{itemize}
  \item Crea uno con \ui{Nuevo estante} en la biblioteca, o desde \ui{Añadir a un estante} en
        cualquier libro.
  \item Haz clic derecho (o mantén pulsado) un estante para renombrarlo o eliminarlo. Al eliminar
        un estante, sus libros siguen en tu biblioteca.
\end{itemize}

\subsection{El menú de un libro}
Haz clic derecho en un libro (o mantenlo pulsado en el teléfono) para ver \ui{Detalles del libro},
\ui{Añadir a un estante}, \ui{Quiero leerlo}, \ui{Marcar como terminado} y dos formas de quitarlo:
\begin{itemize}
  \item \ui{Quitar la descarga} libera espacio, pero deja el libro en tu biblioteca con sus notas,
        estantes y posición de lectura.
  \item \ui{Quitar de la biblioteca} lo quita todo, incluidas las notas y los marcadores.
\end{itemize}

% ---------------------------------------------------------------------------
\section{Leer}

Toca un libro descargado para abrirlo. Se abre donde lo dejaste.

\shot{reader}{El lector con su barra de herramientas.}

\subsection{Pasar páginas}
\begin{keytable}{Acción}{Cómo}
  Página siguiente   & Toca el borde derecho · \kright{} · \kbd{Av Pág} · \kbd{Espacio} \\
  Página anterior    & Toca el borde izquierdo · \kleft{} · \kbd{Re Pág} \\
  Desplazarse un poco & \kup{} \kdown{} o la rueda del ratón \\
  Mostrar la barra   & Toca el centro de la página \\
\end{keytable}

\subsection{La barra de herramientas}
Toca el centro de la página para mostrarla. Arriba: volver a la biblioteca, buscar
(sección~\ref{sec:search}), marcador, índice y ajustes de pantalla. Abajo: capítulo anterior y
siguiente, un control deslizante para saltar entre capítulos y tu progreso.

\subsection{Progreso y tiempo restante}
La línea inferior muestra cuántos minutos quedan del capítulo y cuánto llevas del libro. El tiempo
se calcula con tu propia velocidad de lectura, que la aplicación aprende mientras lees.

\subsection{Índice e imágenes}
\ui{Índice} muestra los capítulos reales del libro, aunque un mismo archivo contenga varios. Las
imágenes aparecen en el texto; tócalas para verlas a pantalla completa y ampliarlas.

% ---------------------------------------------------------------------------
\section{Subrayados, notas y marcadores}

\subsection{Seleccionar texto}
En el teléfono, mantén pulsada una palabra y arrastra los tiradores. En una computadora, arrastra
sobre el texto o haz doble clic en una palabra y luego \textbf{clic derecho} sobre la selección.
Aparece un menú con:
\begin{itemize}
  \item \ui{Subrayar}: marca el texto con tu color de subrayado.
  \item \ui{Nota}: un subrayado con tu propia nota, y a elegir entre cuatro colores.
  \item \ui{Definir} y \ui{Traducir} (sección~\ref{sec:lookup}).
  \item \ui{Copiar}.
\end{itemize}
Los subrayados se ajustan a palabras completas. Los que tienen nota aparecen con una línea debajo.
Vuelve a seleccionar un pasaje subrayado para editar su nota o quitarlo.

\subsection{Marcadores}
Usa el botón de marcador de la barra de herramientas, o \kbd{Ctrl+B}. Una cinta roja indica que
la página tiene un marcador.

\subsection{El panel lateral}
El botón de índice abre un panel con tres pestañas: \ui{Índice}, \ui{Marcadores} y \ui{Notas}. Toca
cualquier elemento para ir a él. \ui{Copiar todo en Markdown} copia todos tus subrayados y notas,
listos para pegarlos en otro programa.

% ---------------------------------------------------------------------------
\section{Diccionario y traducción}\label{sec:lookup}

Selecciona una palabra (hasta tres) y elige \ui{Definir} para ver sus definiciones. Primero
aparecen las del idioma del libro. Provienen de Wiktionary y están escritas en inglés.

Selecciona cualquier pasaje y elige \ui{Traducir}. Elige el idioma de destino en la lista; la
aplicación lo recuerda.

\begin{note}{Se necesita internet}
Las definiciones y traducciones necesitan conexión a internet. El texto seleccionado se envía a
Wiktionary o a MyMemory; no se envía nada más.
\end{note}

% ---------------------------------------------------------------------------
\section{Buscar dentro de un libro}\label{sec:search}

Pulsa \kbd{Ctrl+F} o el botón de búsqueda de la barra de herramientas y escribe al menos dos
letras. La búsqueda no distingue mayúsculas ni acentos. Elige un resultado para ir a él. Una barra
arriba muestra \emph{3 de 27}: usa sus flechas, o \kbd{F3}, para ir al siguiente resultado.
\kbd{Esc} la cierra.

% ---------------------------------------------------------------------------
\section{Pantalla y accesibilidad}\label{sec:display}

Ábrelo con el botón de ajustes de la pantalla de inicio, o con el botón de texto del lector. Una
vista previa arriba muestra cada cambio al momento.

\shot{settings}{Pantalla y accesibilidad.}

\subsection{Idioma}
\ui{Idioma de la aplicación} cambia la interfaz a español, inglés, francés, alemán, italiano o
portugués al instante. \ui{Igual que el dispositivo} sigue el idioma de tu dispositivo.

\subsection{Texto}
\begin{itemize}
  \item Tema: \ui{Claro}, \ui{Sepia} u \ui{Oscuro}.
  \item Tipografía: cuatro tipografías incluidas en la aplicación, para que los libros se vean
        igual en todos los dispositivos:
  \begin{itemize}
    \item \textbf{Literata}, diseñada para leer en pantalla (la predeterminada);
    \item {\garamondfont EB Garamond}, un tipo de libro clásico;
    \item {\sffamily Atkinson Hyperlegible}, creada para que las letras parecidas se distingan
          con facilidad;
    \item {\dyslexicfont OpenDyslexic}, pensada para lectores con dislexia.
  \end{itemize}
  \item \ui{Tamaño del texto}, \ui{Interlineado} y \ui{Espaciado entre letras}.
  \item \ui{Justificar el texto}: desactívalo para que los espacios entre palabras sean iguales,
        algo que a muchos lectores con dislexia les resulta más fácil.
\end{itemize}

\subsection{Página}
\ui{Márgenes} y \ui{Ancho de línea}. Las líneas más cortas se siguen con más facilidad.

\subsection{Pantalla}
\ui{Brillo} atenúa toda la aplicación. \ui{Calidez} reduce la luz azul, como un modo nocturno.
Ambos se suman a los ajustes de tu dispositivo.

\ui{Restablecer valores predeterminados} devuelve todos estos ajustes a su estado inicial.

% ---------------------------------------------------------------------------
\section{Atajos de teclado}

\begin{keytable}{Teclas}{Acción}
  \kright{} \kbd{Av Pág} \kbd{Espacio} & Página siguiente \\
  \kleft{} \kbd{Re Pág}                & Página anterior \\
  \kup{} \kdown{}                      & Desplazarse un poco \\
  \kbd{Ctrl+F}                         & Buscar en el libro \\
  \kbd{F3}                             & Resultado siguiente \\
  \kbd{Ctrl+B}                         & Poner o quitar un marcador en la página \\
  \kbd{Esc}                            & Cerrar la búsqueda o la barra, o volver a la biblioteca \\
  \kbd{Mayús} + flechas                & Ampliar una selección de texto \\
\end{keytable}

% ---------------------------------------------------------------------------
\section{Privacidad y tus datos}

Illuminati no tiene cuentas ni seguimiento. Tu biblioteca, tus notas y tus ajustes se guardan solo
en tu dispositivo (en Windows, en \texttt{\%APPDATA\%\textbackslash Illuminati}).

La aplicación solo se conecta a internet para:
\begin{itemize}
  \item explorar y buscar en los catálogos y descargar libros y portadas;
  \item definir o traducir el texto que seleccionaste, que se envía a Wiktionary o a MyMemory.
\end{itemize}

% ---------------------------------------------------------------------------
\section{Sobre los derechos de autor}

Todos los libros son de dominio público en Estados Unidos. Algunos textos de Wikisource tienen en
cambio licencias libres (como CC BY-SA), que permiten leerlos y compartirlos libremente. De
Internet Archive solo se muestran libros publicados hace 95 años o más.

\begin{tip}{Consulta las leyes de tu país}
Las leyes de derechos de autor varían según el país. Un libro libre en Estados Unidos puede seguir
protegido donde vives, así que consulta las leyes de tu país antes de descargar.
\end{tip}

% ---------------------------------------------------------------------------
\section{Solución de problemas}

\begin{keytable}{Mensaje}{Qué hacer}
  No se pudo conectar con un catálogo & El servicio puede estar ocupado o tu conexión caída.
                             Usa \ui{Reintentar} más tarde; los demás catálogos siguen
                             funcionando. \\
  Este título no tiene edición electrónica & Probablemente es un audiolibro. Busca otra
                             edición. \\
  Wikisource no pudo preparar este libro & La obra es demasiado larga para el servicio de
                             Wikisource. Prueba con uno de sus volúmenes o con otro catálogo. \\
  Se alcanzó el límite diario de traducciones gratuitas & MyMemory permite un número limitado de
                             traducciones gratuitas al día. Inténtalo mañana. \\
  La descarga tardó demasiado & Revisa tu conexión e inténtalo de nuevo. \\
\end{keytable}

% ---------------------------------------------------------------------------
\section{Créditos}

Libros: Project Gutenberg, Standard Ebooks, Wikisource (mediante WS Export) e Internet Archive.
Definiciones: Wiktionary (CC BY-SA). Traducciones: MyMemory.

Tipografías, con licencia SIL Open Font License: Literata, EB Garamond, Atkinson Hyperlegible
(Braille Institute) y OpenDyslexic (Abbie Gonzalez). Sus licencias están en \ui{Acerca de} →
\ui{Ver licencias} dentro de la aplicación.

\end{document}
